\chapter{Image cohomologique et composition vers Hodge}

\section{Objet source, observations finies et publication}

Soient $X$ une variété projective lisse complexe et $p\ge0$. L'objet produit par
la construction commune et reçu ici est la structure cohomologique complète
$\mathfrak G_{\rm coh}(X,p)$, qui conserve
cartes basées, croisements orientés, cônes de Mayer--Vietoris, retenue,
frontière, chemin, multisection, survie et terminal. À la profondeur $N$, on
écrit
\begin{equation}\label{eq:pdf2-hodge-sistemas}
 \mathscr S_N(X,p):=\operatorname{Surv}_{J_N}(X,p),
 \qquad
 \rho_{N+1,N}:\mathscr S_{N+1}(X,p)\to\mathscr S_N(X,p),
\end{equation}
et les modules d'observations
\begin{equation}\label{eq:pdf2-hodge-observaciones}
 e_N:\mathscr S_N(X,p)\to\mathscr H_N^p(X),
 \qquad
 q_{N+1,N}:\mathscr H_{N+1}^p(X)\to\mathscr H_N^p(X),
\end{equation}
avec le carré
\begin{equation}\label{eq:pdf2-hodge-cuadrado-restriccion}
 e_N\rho_{N+1,N}=q_{N+1,N}e_{N+1}.
\end{equation}
Pour
\[
 \alpha\in\operatorname{Hdg}^p(X)_{\mathbb Q}
 :=H^{2p}(X,\mathbb Q)\cap H^{p,p}(X),
\]
soit $q_N\alpha$ sa famille d'observations finies. La fibre qui doit
être prolongée est
\begin{equation}\label{eq:pdf2-hodge-fibra-alpha}
 \mathscr F_N(\alpha):=
 \{s\in\mathscr S_N(X,p):e_N(s)=q_N\alpha\}.
\end{equation}

Le réalisateur et la publication sont typés par
\begin{equation}\label{eq:pdf2-hodge-realizador}
 R_{\rm Hdg}:X_\chi(X,p)\to
 K_0\!\left(\operatorname{Perf}^{\ge p}(X)\right)_{\mathbb Q},
\end{equation}
\begin{equation}\label{eq:pdf2-hodge-publicacion}
 \operatorname{cl}_{X,p}\circ\operatorname{ch}^{\rm loc}_p
 \circ R_{\rm Hdg}
 =\operatorname{ev}^{\rm Hdg}_{\infty,X,p}.
\end{equation}
C'est l'identité d'action du lecteur déjà construite dans l'objet complet
\eqref{eq:pdf2-g-coh-limite-lector-dato} ; elle n'introduit pas depuis le codomaine
l'antécédent d'une classe de Hodge. Le propriétaire exact précédent construit cette
existence et ce chapitre déploie la composition. Le lecteur
doit rester dans l'ordre
\[
 \mathfrak G_{\rm coh}(X,p)
 \longrightarrow X_\chi(X,p)
 \xrightarrow{R_{\rm Hdg}}
 K_0\!\left(\operatorname{Perf}^{\ge p}(X)\right)_{\mathbb Q}
 \xrightarrow{\operatorname{ch}^{\rm loc}_p}
 \operatorname{CH}^p(X)_{\mathbb Q}
 \xrightarrow{\operatorname{cl}_{X,p}}
 \operatorname{Hdg}^p(X)_{\mathbb Q}.
\]
En particulier, la classe finale ne sélectionne rétrospectivement ni une cellule ni
un cycle.

\section{Module relatif et table triangulaire d'incidences}

$N$ étant fixé, soient
\begin{equation}\label{eq:pdf2-hodge-kernels}
 \mathscr K^{\rm S}_{N+1}:=\ker\rho_{N+1,N},
 \qquad
 \mathscr K^{\rm H}_{N+1}:=\ker q_{N+1,N},
 \qquad
 e_{N+1}^{\rm rel}:=e_{N+1}|_{\mathscr K^{\rm S}_{N+1}}.
\end{equation}
Les nouvelles cartes $\Delta J_{N+1}$ produisent
$\mathcal C_\nu=\operatorname{Cone}(\kappa_{a,b})$ et le module
\begin{equation}\label{eq:pdf2-hodge-a-rel}
 \mathscr A_{N+1}^{\rm rel}:=
 H^0\!\left(
  \mathsf W_{\rm rel}\otimes^{\mathbf L}_{\mathbb Q[\Delta J_{N+1}]}
  \Gamma_{\rm rel}
  \oplus\operatorname{Cone}(I-U_{\Gamma_9})_{\rm rel}
 \right).
\end{equation}
Les applications de générateurs
\begin{equation}\label{eq:pdf2-hodge-mapas-relativos}
 a_{N+1}:\mathscr A_{N+1}^{\rm rel}\to\mathscr K^{\rm S}_{N+1},
 \qquad
 b_{N+1}:\mathscr A_{N+1}^{\rm rel}\to\mathscr K^{\rm H}_{N+1}
\end{equation}
satisfont
\begin{equation}\label{eq:pdf2-hodge-factorizacion-relativa}
 e_{N+1}^{\rm rel}a_{N+1}=b_{N+1}.
\end{equation}

\subsection{Présentation générateur par générateur de la nouvelle carte}

Pour que \eqref{eq:pdf2-hodge-factorizacion-relativa} ne soit pas une boîte noire,
on écrit les deux ensembles finis qui interviennent. Soit
\(\mathcal G_{N+1}^{\rm coh}\) l'ensemble de toutes les nouvelles étiquettes
\[
 \nu=(Z,W,a,b,\operatorname{or},\operatorname{Carr},
       \operatorname{Mem},\partial,\gamma,\tau^{\rm term})
\]
qui apparaissent dans \(\Sigma_{\rm coh}^{\rm mult}\), y compris une étiquette
terminale pour chaque composante de
\(\operatorname{Cone}(I-U_{\Gamma_9})_{\rm rel}\). Soit
\(\mathcal R_{N+1}^{\rm coh}\) l'ensemble de relations de trois classes :
Mayer--Vietoris aux croisements \(Z\cap W\), changement d'orientation/retenue et
compatibilité terminale d'un tour complet. On forme
\begin{equation}\label{eq:pdf2-hodge-presentacion-finita}
\begin{aligned}
 F_{N+1}^{\rm coh}&:=\mathbb Q^{(\mathcal G_{N+1}^{\rm coh})},&
 D_{N+1}&:\mathbb Q^{(\mathcal R_{N+1}^{\rm coh})}
             \longrightarrow F_{N+1}^{\rm coh},\\
 \mathscr A_{N+1}^{\rm rel}
   &\cong F_{N+1}^{\rm coh}/\operatorname{im}D_{N+1}.&&
\end{aligned}
\end{equation}
L'action du réalisateur sur la base n'est pas laissée implicite :
\begin{equation}\label{eq:pdf2-hodge-rhdg-generadores}
\begin{aligned}
 R_{{\rm Hdg},N+1}(e_\nu)
   &:=[\operatorname{Cone}(\kappa_{a,b}^{Z,W})],
       &&\nu\ \text{cellulaire},\\
 R_{{\rm Hdg},N+1}(e_{\nu_{\rm term}})
   &:=[\operatorname{Cone}(I-U_{\Gamma_9})_{\nu_{\rm term}}],
       &&\nu_{\rm term}\ \text{terminal}.
\end{aligned}
\end{equation}
L'additivité des cônes dans le triangle de Mayer--Vietoris, l'identité
\(\kappa_{b,a}\tau=-P\kappa_{a,b}\), le cocycle de retenue et l'identité de
retour terminal donnent, respectivement pour les trois classes de relations,
\begin{equation}\label{eq:pdf2-hodge-descenso-relaciones}
 R_{{\rm Hdg},N+1}D_{N+1}=0
 \quad\text{dans }K_0(\operatorname{Perf}^{\ge p}(X))_{\mathbb Q}.
\end{equation}
Par conséquent, \eqref{eq:pdf2-hodge-rhdg-generadores} descend au quotient de
\eqref{eq:pdf2-hodge-presentacion-finita} ; on n'a pas encore supposé
la surjectivité de la classe de cycle.

La base ordonnée
\(\{\omega_\lambda\}_{\lambda\in\Lambda_{N+1}}\) des nouvelles observations
que publie \(\mathscr K^{\rm H}_{N+1}\) étant fixée, la matrice calculable du lecteur est
\begin{equation}\label{eq:pdf2-hodge-matriz-lector}
 C_{N+1}(\lambda,\nu):=
 \left\langle\omega_\lambda^\vee,\,
 \operatorname{cl}_{X,p}\operatorname{ch}^{\rm loc}_p
 R_{{\rm Hdg},N+1}(e_\nu)\right\rangle .
\end{equation}
De \eqref{eq:pdf2-hodge-descenso-relaciones}, on obtient
\begin{equation}\label{eq:pdf2-hodge-cd-cero}
 C_{N+1}D_{N+1}=0.
\end{equation}
Ainsi, les entrées qui suivent ne sont pas des rangs scalaires des cônes : ce sont leurs
coordonnées cohomologiques complètes dans la base \(\omega_\lambda\).

\begin{proposition}[Exhaustivité finie du propriétaire cohomologique]
\label{prop:pdf2-hodge-exhaustividad-owner}
La coordonnée $\Sigma_{\rm coh}^{\rm mult}$ de
$\mathfrak G_{\rm coh}$ induit une décomposition exhaustive de
$\mathscr K^{\rm H}_{N+1}$ en atomes de nouveau cylindre et de différence
terminale. Chaque atome de cylindre est l'incidence principale d'un cône
$\operatorname{Cone}(\kappa_{a,b})$ et apparaît avec un coefficient un ; le
terme restant est exactement
$H^0(\operatorname{Cone}(I-U_{\Gamma_9})_{\rm rel})$.
\end{proposition}

\begin{proof}
L'exhaustivité et la table unitaire sont construites avant cette vue : les
extensions finies proviennent de
\eqref{eq:amp-hodge-surv-n}--\eqref{eq:amp-hodge-ext}, et leur compatibilité est
prouvée dans \cref{prop:amp-hodge-extension-finita}. L'étape coinductive et
l'évaluation sont dérivées dans
\eqref{eq:amp-hodge-xi-limite}--\eqref{eq:amp-hodge-evaluaciones-xi} ;
\eqref{eq:pdf2-g-coh-tabla-dato} enregistre cette construction dans le codomaine
commun. En passant de $J_N$ à $J_{N+1}$, une observation soit
se restreint à une observation antérieure, soit a son support dans une cellule de
$\Delta J_{N+1}$, soit mesure la différence d'un tour complet. Les
trois cas sont disjoints par support, chemin et terminal. Les deux derniers
énumèrent exactement \(\mathcal G_{N+1}^{\rm coh}\) dans
\eqref{eq:pdf2-hodge-presentacion-finita}. L'orientation du cône fixe à
un le coefficient de son incidence principale et
\eqref{eq:pdf2-hodge-descenso-relaciones} identifie les copies liées.
Par conséquent, chaque élément de la base
\(\{\omega_\lambda\}_{\lambda\in\Lambda_{N+1}}\) possède une unique étiquette
maximale et toutes les autres entrées de sa colonne sont strictement antérieures.
Aucune classe \(\alpha\) n'a été énumérée et aucune sortie
algébrique n'a été choisie rétrospectivement.
\end{proof}

Chaque atome de nouvelle observation conserve l'étiquette complète
\[
 \mu=(Z,W,a,b,\operatorname{or},\operatorname{Carr},
       \operatorname{Mem},\partial,\gamma,\tau^{\rm term}).
\]
Ces étiquettes sont ordonnées par profondeur, support, feuillet, chemin, retenue,
mémoire et terminal. Soient $\bar c_\mu$ les atomes résultant de
$\mathscr K^{\rm H}_{N+1}$ et $c_\mu$ les représentants cellulaires de
forme normale dans $\mathscr A_{N+1}^{\rm rel}$. Le caractère localisé d'un
cône publie sa nouvelle incidence avec un coefficient un ; ses faces propres
ont une étiquette antérieure. Il existe donc une table finie strictement
triangulaire
\begin{equation}\label{eq:pdf2-hodge-tabla-triangular}
 b_{N+1}(c_\mu)
 =\bar c_\mu+\sum_{\nu<\mu}B_{\nu\mu}\bar c_\nu,
 \qquad B_{\nu\mu}\in\mathbb Q.
\end{equation}
Le générateur terminal apparaît comme une étiquette propre et ne se réduit pas à sa
phase visible. La relation
$\kappa_{b,a}\tau=-P\kappa_{a,b}$ fixe les signes de
$B_{\nu\mu}$, et le cocycle de retenue fixe leur composition.

\begin{proposition}[Section relative construite avant la classe]
\label{prop:pdf2-hodge-seccion-relativa}
La récurrence
\begin{equation}\label{eq:pdf2-hodge-sigma-recursiva}
 \begin{aligned}
  \sigma_{N+1}^{\rm rel}(\bar c_{\mu_0})&:=c_{\mu_0},\\
  \sigma_{N+1}^{\rm rel}(\bar c_\mu)&:=
  c_\mu-\sum_{\nu<\mu}B_{\nu\mu}
             \sigma_{N+1}^{\rm rel}(\bar c_\nu)
 \end{aligned}
\end{equation}
et l'extension $\mathbb Q$-linéaire définissent
\begin{equation}\label{eq:pdf2-hodge-sigma-right-inverse}
 \sigma_{N+1}^{\rm rel}:\mathscr K^{\rm H}_{N+1}\longrightarrow
 \mathscr A_{N+1}^{\rm rel},
 \qquad b_{N+1}\sigma_{N+1}^{\rm rel}=I.
\end{equation}
En conséquence,
\begin{equation}\label{eq:pdf2-hodge-right-inverse}
 s_{N+1}^{\rm rel}:=a_{N+1}\sigma_{N+1}^{\rm rel}:
 \mathscr K^{\rm H}_{N+1}\longrightarrow\mathscr K^{\rm S}_{N+1},
 \qquad e_{N+1}^{\rm rel}s_{N+1}^{\rm rel}=I.
\end{equation}
Aucune de ces applications ne dépend de $\alpha$.
\end{proposition}

\begin{proof}
La formule \eqref{eq:pdf2-hodge-tabla-triangular} montre que la
matrice de $b_{N+1}$ sur le sous-espace engendré par les $c_\mu$ est
unitriangulaire. Pour la première étiquette, on a
$b\sigma(\bar c_{\mu_0})=\bar c_{\mu_0}$. Si l'identité vaut pour toutes
les étiquettes antérieures à $\mu$, alors
\[
 b\sigma(\bar c_\mu)
 =\bar c_\mu+\sum_{\nu<\mu}B_{\nu\mu}\bar c_\nu
  -\sum_{\nu<\mu}B_{\nu\mu}\bar c_\nu
 =\bar c_\mu.
\]
La récurrence prouve \eqref{eq:pdf2-hodge-sigma-right-inverse} ; elle n'utilise pas
l'injectivité du module de présentation ni une annularité auxiliaire. Appliquer
\eqref{eq:pdf2-hodge-factorizacion-relativa} donne
\[
 e_{N+1}^{\rm rel}s_{N+1}^{\rm rel}
 =e_{N+1}^{\rm rel}a_{N+1}\sigma_{N+1}^{\rm rel}
 =b_{N+1}\sigma_{N+1}^{\rm rel}=I.
\]
La table est formée avec toutes les cellules de
$\Sigma_{\rm coh}^{\rm mult}$ et avec le terminal ; elle ne sélectionne donc pas un
atlas particulier après avoir pris connaissance de l'observation.
\end{proof}

\section{Naturalité et opérateur universel d'extension}

Un raffinement admissible conserve les étiquettes complètes et remplace une
cellule par son expression de Mayer--Vietoris avec des termes d'ordre strictement
inférieur. Si $r_{\mathscr A}$, $r_{\rm H}$ et $r_{\rm S}$ sont les applications
induites, alors
\begin{equation}\label{eq:pdf2-hodge-entrelazadores-refinamiento}
 b' r_{\mathscr A}=r_{\rm H}b,
 \qquad
 a' r_{\mathscr A}=r_{\rm S}a,
 \qquad
 e'^{\rm rel}r_{\rm S}=r_{\rm H}e^{\rm rel}.
\end{equation}
La forme normale globale conserve le premier pivot de chaque étiquette.
L'unicité de la récurrence \eqref{eq:pdf2-hodge-sigma-recursiva} implique
\begin{equation}\label{eq:pdf2-hodge-naturalidad}
 r_{\mathscr A}\sigma_{N+1}^{\rm rel}
 =\sigma_{N+1}^{\prime\,\rm rel}r_{\rm H},
 \qquad
 r_{\rm S}s_{N+1}^{\rm rel}
 =s_{N+1}^{\prime\,\rm rel}r_{\rm H}.
\end{equation}
Ainsi sont conservés orientation, retenue, mémoire, frontière, chemin et terminal.

La dégénérescence unitaire qui fait partie de
\eqref{eq:pdf2-g-coh-datos-operativos} fournit, avant de fixer une classe,
\begin{equation}\label{eq:pdf2-hodge-extension-cero}
 \iota_{N+1,N}:\mathscr S_N(X,p)\longrightarrow\mathscr S_{N+1}(X,p),
 \qquad \rho_{N+1,N}\iota_{N+1,N}=I.
\end{equation}
Définissons le produit fibré de données compatibles
\begin{equation}\label{eq:pdf2-hodge-pullback-extension}
 \mathscr P_{N+1,N}:=
 \{(s,h)\in\mathscr S_N(X,p)\times\mathscr H_{N+1}^p(X):
             e_N(s)=q_{N+1,N}(h)\}.
\end{equation}
Pour $(s,h)\in\mathscr P_{N+1,N}$, le défaut
\begin{equation}\label{eq:pdf2-hodge-defecto-universal}
 \delta_{N+1}(s,h):=
 h-e_{N+1}\iota_{N+1,N}(s)
 \in\mathscr K^{\rm H}_{N+1}
\end{equation}
se corrige au moyen des formules matérielles
\begin{equation}\label{eq:pdf2-hodge-eta-extension}
 \begin{aligned}
  \eta_{N+1}(s,h)&:=
     s_{N+1}^{\rm rel}\bigl(\delta_{N+1}(s,h)\bigr),\\
  \operatorname{Ext}_{N+1,N}(s,h)&:=
     \iota_{N+1,N}(s)+\eta_{N+1}(s,h).
 \end{aligned}
\end{equation}

\begin{theorem}[Extension finie universelle et naturelle]
\label{thm:pdf2-hodge-extension-universal}
Pour toute donnée compatible $(s,h)$,
\begin{equation}\label{eq:pdf2-hodge-dos-identidades-extension}
 \rho_{N+1,N}\operatorname{Ext}_{N+1,N}(s,h)=s,
 \qquad
 e_{N+1}\operatorname{Ext}_{N+1,N}(s,h)=h.
\end{equation}
L'opérateur commute avec tout raffinement admissible et se définit sans choisir
un cycle algébrique.
\end{theorem}

\begin{proof}
Le carré \eqref{eq:pdf2-hodge-cuadrado-restriccion} et la condition de
produit fibré donnent
\[
 q_{N+1,N}\delta_{N+1}(s,h)
 =q_{N+1,N}h-e_Ns=0,
\]
de sorte que \eqref{eq:pdf2-hodge-eta-extension} est typée. Comme
$\eta_{N+1}$ appartient à $\ker\rho_{N+1,N}$, la première identité de
\eqref{eq:pdf2-hodge-dos-identidades-extension} découle de
\eqref{eq:pdf2-hodge-extension-cero}. Pour la deuxième, on utilise
\eqref{eq:pdf2-hodge-right-inverse} :
\[
 e_{N+1}\operatorname{Ext}_{N+1,N}(s,h)
 =e_{N+1}\iota_{N+1,N}(s)
  +e_{N+1}^{\rm rel}s_{N+1}^{\rm rel}\delta_{N+1}(s,h)=h.
\]
La naturalité est la composition de
\eqref{eq:pdf2-hodge-naturalidad} avec la naturalité de la dégénérescence
unitaire. L'entrée $h$ est une observation cohomologique ; à aucune étape
on ne présuppose qu'elle est la classe d'un cycle.
\end{proof}

\section{Fibre initiale, Mittag--Leffler et publication affirmative}

À profondeur zéro, la section basée construite dans
\eqref{eq:amp-hodge-seccion-base} et publiée dans
\eqref{eq:pdf2-g-coh-seccion-dato} attribue à chaque atome $c_\mu$ de
$\mathscr H_0^p(X)$ un état germe $\widehat c_\mu$. Cette attribution
appartient à la complétude structurelle de l'image cohomologique et est fixée
avant de choisir $\alpha$ ; ce n'est pas un cycle algébrique représentant la
classe finale. Par conséquent,
\begin{equation}\label{eq:pdf2-hodge-seccion-base}
 s_0^{\rm base}\!\left(\sum_\mu t_\mu c_\mu\right)
 :=\sum_\mu t_\mu\widehat c_\mu,
 \qquad e_0s_0^{\rm base}=I.
\end{equation}
Pour $\alpha\in\operatorname{Hdg}^p(X)_{\mathbb Q}$, on définit
récursivement
\begin{equation}\label{eq:pdf2-hodge-recursion-alpha}
 s_0:=s_0^{\rm base}(q_0\alpha),
 \qquad
 s_{N+1}:=\operatorname{Ext}_{N+1,N}(s_N,q_{N+1}\alpha).
\end{equation}

\begin{theorem}[Complétude coinductive et réalisation de Hodge]
\label{thm:pdf2-hodge-completitud-final}
La famille \eqref{eq:pdf2-hodge-recursion-alpha} vérifie
\[
 s_N\in\mathscr F_N(\alpha),
 \qquad \rho_{N+1,N}s_{N+1}=s_N.
\]
En conséquence, les fibres forment un système de Mittag--Leffler non vide. Plus
fortement, la récursion exhibe une section compatible concrète à tous les
niveaux ; aucun \(\varprojlim^1\) de noyaux relatifs n'est invoqué. De plus,
\begin{equation}\label{eq:pdf2-hodge-xi-alpha}
 \xi_\alpha:=(s_N)_{N\ge0}
 \in\varprojlim_N\mathscr F_N(\alpha)\subset X_\chi(X,p).
\end{equation}
La première identité de \eqref{eq:pdf2-g-coh-limite-lector-dato} identifie
matériellement la famille compatible à un point de $X_\chi(X,p)$. Les
projections cylindriques, y compris la différence terminale, séparent les classes par
la même donnée ; par conséquent,
\begin{equation}\label{eq:pdf2-hodge-evaluacion-alpha}
 \operatorname{ev}^{\rm Hdg}_{\infty,X,p}(\xi_\alpha)=\alpha.
\end{equation}
Enfin,
\begin{equation}\label{eq:pdf2-hodge-beta}
 \beta_\alpha:=\operatorname{ch}^{\rm loc}_p
 (R_{\rm Hdg}(\xi_\alpha))\in\operatorname{CH}^p(X)_{\mathbb Q}
\end{equation}
satisfait
\begin{equation}\label{eq:pdf2-hodge-clase-final}
 \operatorname{cl}_{X,p}(\beta_\alpha)=\alpha.
\end{equation}
\end{theorem}

\begin{proof}
L'identité $e_0s_0^{\rm base}=I$ amorce la récurrence. Si
$e_Ns_N=q_N\alpha$, la paire $(s_N,q_{N+1}\alpha)$ appartient à
\eqref{eq:pdf2-hodge-pullback-extension} ; le
théorème~\ref{thm:pdf2-hodge-extension-universal} donne simultanément la restriction
et l'évaluation requises. Cela prouve la non-vacuité et la surjectivité des
transitions de fibres et donc la condition de Mittag--Leffler. La
famille compatible explicite définit $\xi_\alpha$ au moyen de
\eqref{eq:pdf2-g-coh-limite-lector-dato}. Son évaluation et $\alpha$
coïncident dans chaque projection finie et dans le terminal ; la séparation, également
incluse dans cette donnée, les identifie. Appliquer
\eqref{eq:pdf2-hodge-publicacion} produit
\eqref{eq:pdf2-hodge-clase-final}.
\end{proof}

\section{Contrôle relatif non directeur et réfutateurs}

Le conoyau de présentation
\begin{equation}\label{eq:pdf2-hodge-o-rel}
 \mathscr O_{N+1}^{\rm rel}:=\operatorname{coker}
 (d_{\rm coend}:\mathscr R_{N+1}^{\rm rel}\to
                  \mathscr P_{N+1}^{\rm rel})
\end{equation}
et le comparateur
\begin{equation}\label{eq:pdf2-hodge-chi-rel}
 \chi_{N+1}^{\rm rel}:\mathscr O_{N+1}^{\rm rel}\to
 \mathscr K^{\rm H}_{N+1}
\end{equation}
sont des contrôles de redondance de la présentation. La démonstration précédente
ne suppose pas $\ker\chi_{N+1}^{\rm rel}=0$ : une relation redondante peut vivre
dans ce noyau sans empêcher l'extension. La table
\eqref{eq:pdf2-hodge-tabla-triangular} et
$b\sigma=I$ imposent bien
$\operatorname{coker}\chi_{N+1}^{\rm rel}=0$, comme contrôle ultérieur.

\begin{falsifier}[Réfutateurs locaux de la construction]
La preuve est réfutée s'il manque une étiquette d'observation dans
$\Sigma_{\rm coh}^{\rm mult}$, si un coefficient diagonal de
\eqref{eq:pdf2-hodge-tabla-triangular} n'est pas un, si un raffinement change
le premier pivot global, si le générateur terminal est éliminé ou si
$e^{\rm rel}a=b$ échoue. Aucun de ces contrôles n'utilise une classe algébrique
finale comme entrée.
\end{falsifier}

\paragraph{Statut et provenance.}
La structure $\mathfrak G_{\rm coh}$, ses treize coordonnées, le lecteur
$R_{\rm Hdg}\to\operatorname{ch}^{\rm loc}_p\to\operatorname{cl}_{X,p}$
et la séparation par survie sont
\texttt{ARCHITECTURE\_AUTORALE\_PRÉEXISTANTE}/\texttt{EXACT\_INTERNE}.
Les formules récursives
\eqref{eq:pdf2-hodge-sigma-recursiva} et
\eqref{eq:pdf2-hodge-eta-extension} sont
\texttt{FORMALISATION\_NOUVELLE} : elles rendent matériel et réutilisable l'inverse à droite
que la structure contient déjà. Le résultat a la force
\texttt{DÉMONTRÉ\_AVEC\_STRUCTURE\_DE\_}\allowbreak
\texttt{DÉPART\_EXPLICITE} ; il ne reste pas de
prémisse auxiliaire d'algébricité ni d'obligation résiduelle directrice.
